\section*{Chapter \ref{ch:b2:gratings} --- Multiple-Wave Interference and Gratings}

\begin{solution}{exo:b2:gratings:1}
$a = \qty{2}{\micro m}$, $\sin\theta_p = 0.275p$: \qty{16.0}{\degree}, \qty{33.4}{\degree}, \qty{55.6}{\degree};
no fourth order. First-order visible from \qty{11.5}{\degree} to \qty{20.5}{\degree}:
\qty{9}{\degree}. At \qty{30}{\degree}: $\sin\theta = 0.5 + 0.275p$: $p = +1$ (\qty{51}{\degree}) on one
side, $p = -1$ to $-5$ ($13^\circ$, $-3^\circ$, $-19^\circ$, $-37^\circ$, $-61^\circ$) on the
other.
\end{solution}

\begin{solution}{exo:b2:gratings:2}
$a = \qty{1.67}{\micro m}$. Order 1: $\theta = \qty{20.7}{\degree}$, $\Delta\theta = \Delta\lambda/a\cos\theta =
\qty{3.9e-4}{rad}$, \qty{0.19}{mm}; order 2: $\theta = \qty{45}{\degree}$, $\qty{1.0e-3}{rad}$,
\qty{0.51}{mm}. Pixels of \qty{50}{\micro m} or less.
\end{solution}

\begin{solution}{exo:b2:gratings:3}
$N = 12\,000$: $R = 12\,000$, $24\,000$, $36\,000$; $\Delta\lambda = 0.046$, $0.023$, \qty{0.015}{nm};
the doublet (\qty{0.6}{nm}) and the H/D pair (\qty{0.18}{nm}, needing $R = 3600$)
are easily separated.
\end{solution}

\begin{solution}{exo:b2:gratings:4}
CD: violet \qty{14.5}{\degree}, red \qty{25.9}{\degree}; red exists to order 2, violet to
order 4: two (partly overlapping) rainbows. DVD: violet \qty{32.7}{\degree}, red
\qty{71}{\degree}; red has only the first order: one rainbow, wider.
\end{solution}

\begin{solution}{exo:b2:gratings:5}
(a) Zeros at $\varphi = 2\pi m/N$, $m = 1, \dots, N - 1$: $N - 1$ zeros, $N - 2$ maxima
between them. (b) At $N\varphi/2 = 3\pi/2$, $I \approx I_0/\sin^2(3\pi/2N) \approx I_0(2N/3\pi)^2$:
$4/9\pi^2 = 0.045$ of $N^2I_0$. (c) $N = 3$: a single \omterm{thm:b2:gratings:nwaves}{secondary maximum} at
$\varphi = \pi$ of height $I_0$, one ninth of the principal $9I_0$; $N = 6$: four
secondary maxima of a few percent. (d) The central lobe of $\sin^2x/x^2$
holds $90\%$ of the energy.
\end{solution}

\begin{solution}{exo:b2:gratings:6}
(a) $p > \lambda_{\min}/(\lambda_{\max} - \lambda_{\min}) = 1.33$: from $p = 2$. (b) $\lambda/p$. (c)
\qty{25}{nm}; a prism crossed with the grating stacks the orders one above
the other. (d) $\Delta\lambda = \lambda^2/2L$, a fraction of a nanometre --- hundreds
of times smaller than a grating's.
\end{solution}

\begin{solution}{exo:b2:gratings:7}
(a) $\sin\beta = \lambda/2a = 0.3$, $\beta = \qty{17.5}{\degree}$. (b) The single facet diffracts
into a broad lobe centred on its specular direction; the orders that
fall inside it get the light. (c) Order 2 ($2 \times \qty{500}{nm} = \qty{1}{\micro m}$ at
the same angle). (d) $R = 2W\sin\beta/\lambda = 3.6 \times 10^5$.
\end{solution}

\begin{solution}{exo:b2:gratings:8}
(a) $\dd\lambda/\dd x = a\cos\theta/pf = \qty{3.1}{nm/mm}$: \qty{50}{\micro m} is \qty{0.16}{nm}. (b)
$R = 30\,000$: \qty{0.018}{nm}. (c) The slit, by nine; equal at $s = \qty{6}{\micro m}$.
(d) Diffraction at the slit and a starved detector.
\end{solution}

\begin{solution}{exo:b2:gratings:9}
(a) $p = 2L/\lambda = 16\,700$; FSR $c/2L = \qty{30}{GHz}$, $\lambda^2/2L = \qty{0.036}{nm}$. (b) $\mathcal F
= \pi\sqrt{0.95}/0.05 = 61$; \qty{0.5}{GHz}, \qty{0.6}{pm}. (c) $p\mathcal F = 10^6$, four
times the \qty{10}{cm} grating's. (d) Its \omterm{prop:b2:gratings:fabryperot}{free spectral range} is a fraction
of a nanometre: the orders overlap unless the light is prefiltered
to one of them.
\end{solution}

\begin{solution}{exo:b2:gratings:10}
(a) $|\sin\theta - \sin\theta_0| \le 2$. (b) The extreme rays differ in path by
$W(\sin\theta - \sin\theta_0)$ at most $2W$. (c) $2 \times 0.3/5 \times 10^{-7} = 1.2 \times 10^6$; $c/R
= \qty{250}{m/s}$. (d) A line's centre is located to a hundredth of its
width, and a thousand lines average down by $\sqrt{1000}$: metres per
second, with heroic stability.
\end{solution}

\begin{solution}{exo:b2:gratings:11}
(a) The wave reflected by the lower plane travels $2d\sin\theta$ more.
(b) $\sin\theta = 0.273p$: \qty{15.8}{\degree}, \qty{33.1}{\degree}, \qty{55.0}{\degree}. (c) $\lambda \gg 2d$:
no angle satisfies the equation; glass has no regular planes, hence
only diffuse rings --- it is amorphous. (d) $\sin\theta \le 1$.
\end{solution}

\begin{solution}{exo:b2:gratings:12}
(a) Neighbours differ in phase by $2\pi a\sin\theta/\lambda$ from the path and $-\psi$
from the feed; the \omterm{thm:b2:gratings:nwaves}{principal maximum} at $\varphi = 0$. (b) First zero at
$\Delta\varphi = 2\pi/N$: $\Delta\theta = \lambda/Na\cos\theta$. (c) A second \omterm{thm:b2:gratings:nwaves}{principal maximum} at
$\varphi = \pm2\pi$ needs $|\sin\theta \pm \lambda/a| \le 1$: impossible for every steering
angle only if $a \le \lambda/2$. (d) $Na = \qty{0.96}{m}$: \qty{1.8}{\degree}; $50$ steps,
\qty{50}{\micro s}.
\end{solution}

\begin{solution}{pb:b2:gratings:1}
\textbf{1.} $a = \qty{0.833}{\micro m}$; $N = 144\,000$.

\textbf{2.} $\sin\theta = \sin20^\circ - \lambda/a$ (the order lies on the other side of
the normal): \qty{550}{nm} at \qty{18.5}{\degree}, \qty{400}{nm} at \qty{7.9}{\degree}, \qty{700}{nm} at
\qty{29.9}{\degree}: \qty{22}{\degree} of spectrum.

\textbf{3.} $1/a\cos\theta = \qty{1.27e-3}{rad/nm}$; $f_2\times$: \qty{1.27}{mm/nm}, i.e.\
\qty{0.79}{nm/mm}, \qty{0.012}{nm} per pixel.

\textbf{4.} $R = N = 144\,000$: \qty{0.0038}{nm}, a third of a pixel.

\textbf{5.} Second-order \qty{400}{nm} lands where first-order \qty{800}{nm}
would, beyond the red: no overlap in the visible; the overlap would
start at \qty{350}{nm}, which a glass filter removes.

\textbf{6.} \qty{22}{\degree} $\times$ \qty{1}{m} $= \qty{38}{cm}$: a detector covers a slice;
the grating is turned to choose it.

\textbf{7.} $\lambda = 2a\sin17^\circ = \qty{490}{nm}$.

\textbf{8.} $0.8 \times 0.4 = 32\%$.

\textbf{9.} Angles \qty{26.5}{\degree}, \qty{14.0}{\degree}, \qty{10.3}{\degree}: $+14$, $-8$ and
\qty{-14}{cm} from the \qty{550}{nm} position.

\textbf{10.} \qty{0.18}{nm} is $15$ pixels: separated.

\textbf{11.} $50$ pixels apart, each $4$ pixels wide.

\textbf{12.} \qty{100}{\micro m}, $6.7$ pixels, \qty{0.08}{nm} --- twenty times the
grating's limit: slit-limited.

\textbf{13.} The image is \qty{97}{\micro m} wide: a \qty{50}{\micro m} slit passes
half the light.

\textbf{14.} Less light for the same sharpness: the photon noise grows.

\textbf{15.} $\Delta\lambda = \lambda v/c = \qty{0.066}{nm}$, $5.5$ pixels.

\textbf{16.} The Earth's velocity shifts every line by up to $\pm5.5$
pixels through the year; it is known from the ephemerides to far
better than the $0.3\%$ needed for \qty{100}{m/s}.

\textbf{17.} \qty{2.2e-5}{nm}, $0.002$ pixel --- far below the line width. A
centroid to a hundredth of a \qty{7}{px} line is \qty{0.07}{px}; a thousand
lines give $\sqrt{1000}$: $0.002$ pixel --- the signal, just. Real
instruments add resolving power and stability.

\textbf{18.} $\Delta\sin\theta = (\lambda/a) \times 10^{-5} = 6.6 \times 10^{-6}$, $\Delta\theta = \qty{7e-6}{rad}$,
\qty{7}{\micro m}: half a pixel per kelvin --- millikelvins for \qty{10}{m/s}.

\textbf{19.} Its lines, recorded at the same instant on the same
pixels, give the wavelength scale and track every drift.

\textbf{20.} $a = \qty{3.33}{\micro m}$, $N = 36\,000$, same angle (\qty{18.5}{\degree}), $R = 4N
= 144\,000$, same dispersion, FSR $\lambda/4 = \qty{137}{nm}$: resolving power and
dispersion depend on $W$ and the angle only; the \omterm{prop:b2:gratings:fabryperot}{free spectral range}
shrank.

\textbf{21.} FSR $\lambda^2/2L = \qty{0.15}{nm}$ ($12$ pixels); $\mathcal F = 30$; peaks
\qty{0.005}{nm} wide: a comb of sharp lines across the detector, a
wavelength ruler.

\textbf{22.} $2 \times 10^{-4}$ against \qty{1.3e-3}{rad/nm}: six times less, and a
prism's resolving power ($b\,\dd n/\dd\lambda \sim 10^4$) and non-linear scale
lose too; gratings also work in the ultraviolet.

\textbf{23.} $R \le 2W/\lambda$: $4 \times 10^6$ for a metre at \qty{500}{nm} --- width and
angle are the only ways up.

\textbf{24.} $10^4 \times 0.012 \times 0.32 = 38$ photons per pixel per second; $3800$
in \qty{100}{s}: SNR $\approx 60$.

\textbf{25.} Dispersion, set by $p/a\cos\theta$ and the camera's focal length;
resolving power, by the illuminated width and the angles ($pN$);
\omterm{prop:b2:gratings:fabryperot}{free spectral range}, by the order ($\lambda/p$).
\end{solution}
